There are some limitations in the current implementation that affect the evaluation results. As discussed earlier, our approach performs a widening step from 32-bit to 64-bit values. This helps expose opportunities for extension elimination, but at the same time it can introduce additional extension instructions. Although many of these instructions are removed during the transformation, some of them remain in the generated code.

The main reason for this is that our implementation does not include the backward analysis phase of the original Kawahito et al. algorithm~\cite{kawahito2002effective}, called AnalyzeDEF. This phase is used to find additional extension elimination opportunities and would remove many of the extra trunc-sext instructions that remain after the widening step.

During the development of the pass, we observed that many extension elimination opportunities depend on this backward analysis. For this reason, implementing AnalyzeDEF would likely allow more redundant extension instructions to be removed. It could also create additional opportunities for moving extension operations to less frequently executed regions of the program.

Another possible improvement to the current implementation would be the use of LLVM's Demanded Bits Analysis~\cite{llvm-demandedbits}. This analysis identifies which bits of a value are actually needed by the instructions that use it and could help expose more extension elimination opportunities during the forward analysis phase. By providing more information about how values are used, it could allow the transformation to handle more cases and eliminate a larger number of redundant extension instructions.
